% C/C++ API 兼容性
% license Usr
% type Note


\begin{itemize}
\item 跨 dll/so 调用\textbf{符号（包括函数、全局变量等）}时，C API 是稳定的，意思是统一操作系统和 CPU 架构中，不同的编译器或同一编译器的不同版本编译的 dll/so 等可以混用。 但 C++ API 不稳定，不能混用。
\item 在 C++ 中声明了 \verb`extern "C" {...}` 的部分也是 C API。
\item dll/so 中的符号可以设置编译器属性（非标准 C/C++）如 GCC 的 \verb`__attribute__((attribute_name))` MSVC 和 MinGW 的 \verb`__declspec(attribute_name)`。注意双层括号不可省略为单层。其中 \verb`attribute_name` 在 \verb`__declspec` 中可以是 \verb`dllimport` 或 \verb`dllexport`；若用 GCC/Clang 在非 Windows 系统跨平台使用则用 \verb`visibility("default")` 表示 public，\verb`visibility("hidden")` 表示 hidden（不区分 import/export）。
\item \verb`default/hidden` 是对应了编译器选项 \verb`-fvisibility=default/hidden`。
\item GCC 默认 dll/so 库中的符号全部导出（即 \verb`-fvisibility=default`），而 MSVC 默认全部隐藏。 只有导出的符号才可以被其他库或可执行文件使用。
\item GCC 若要默认隐藏，可以通过（即 \verb`-fvisibility=hidden`）控制某个编译单元，或者用 linker 选项 \verb`-Wl,--exclude-all-symbols` 隐藏当前编译的库的所有符号。
\begin{lstlisting}[language=cpp]
// mylib.h
// MYLIB_BUILDING 宏在编译 mylib.dll 时必须定义，在使用时必须没有定义

#if defined(_MSC_VER)
    // MSVC handling
    #ifdef MYLIB_BUILDING
        #define MYLIB_API __declspec(dllexport)
    #else
        #define MYLIB_API __declspec(dllimport)
    #endif
#elif defined(__GNUC__) || defined(__clang__)
    // GCC/Clang - cross other platforms
    #ifdef MYLIB_BUILDING
        #define MYLIB_API __attribute__((visibility("default")))
    #else
        #define MYLIB_API
    #endif
#else
    #define MYLIB_API  // Unknown compiler
#endif

class MYLIB_API MyClass {  // Exported on all platforms
public:
    void method();
};
\end{lstlisting}
\end{itemize}
